\chapter{Cyclic and permutation groups}\label{chap:M07.02}

\KEntry{K07.02.cyclic_definition}{Cyclic groups}
By Proposition~\ref{prop:integer-powers}, the set
$\langle x\rangle=\{x^k:k\in\Z\}$ is a subgroup containing $x$, and every subgroup containing $x$ contains these powers. Thus it agrees with the generated subgroup of Chapter~\ref{chap:M07.01}. A group is cyclic if it equals $\langle x\rangle$ for some $x$. This definition includes the trivial group.

\KEntry{K07.02.element_order}{Finite and infinite order}
The order $\operatorname{ord}(x)$ is the least positive $n$ satisfying $x^n=1$, if there is one; otherwise its order is infinite. In particular $\operatorname{ord}(x)=1$ if and only if $x=1$, because $x^1=x$. Mathlib's natural-valued \texttt{orderOf} uses $0$ for infinite order. It does not assert a positive annihilating exponent in that case; all statements below distinguish the two cases explicitly.

\KEntry{K07.02.order_divisibility}{When powers agree}
\begin{proposition}\label{prop:order-divisibility}
If $x$ has order $n<\infty$, then for integers $k,a,b$,
\[
 x^k=1\ \Longleftrightarrow\ n\mid k,
 \qquad x^a=x^b\ \Longleftrightarrow\ n\mid(a-b).
\]
If $x$ has infinite order, $x^a=x^b$ if and only if $a=b$.
\end{proposition}
\begin{proof}
Write $k=qn+r$ with $0\leq r<n$ by integer division, valid also for negative $k$. The power laws give $x^k=x^r$. If this is $1$, minimality of $n$ forces $r=0$; the converse follows from $x^n=1$. The second assertion follows by multiplying by $x^{-b}$. For infinite order, an equality with $a\ne b$ would give $x^{a-b}=1$ and, after inverting if necessary, a positive annihilating exponent, a contradiction. Equality for $a=b$ is immediate.
\end{proof}

\KEntry{K07.02.cyclic_models_isomorphism}{The models for cyclic groups}
\begin{theorem}\label{thm:cyclic-models}
A cyclic group generated by an infinite-order element is isomorphic to $(\Z,+)$. A group generated by an element of finite order $n$ is isomorphic to $(\Z/n\Z,+)$ and has $n$ elements. Thus two cyclic groups are isomorphic if and only if they have the same order, distinguishing finite order from infinity.
\end{theorem}
\begin{proof}
The map $k\mapsto x^k$ is a surjective homomorphism from $\Z$ by the power laws. In the infinite case Proposition~\ref{prop:order-divisibility} makes it injective. In the finite case it descends to $[k]\mapsto x^k$; that proposition proves both well-definedness and injectivity on residue classes. Integer division shows there are exactly $n$ such classes. Groups isomorphic to the same model are isomorphic to each other. Conversely an isomorphism is a set bijection, preserves powers, and hence preserves the order of every element and the finite or infinite cardinality of a cyclic group. We write $C_n$ for the finite model, with $C_1$ trivial.
\end{proof}

\KEntry{K07.02.subgroups_of_integer_group}{Subgroups of the integers}
\begin{proposition}\label{prop:integer-subgroups}
Every additive subgroup of $\Z$ is $d\Z$ for a unique $d\geq0$. If $d>0$ it is infinite cyclic; $0\Z=\{0\}$.
\end{proposition}
\begin{proof}
If the subgroup is zero take $d=0$. Otherwise it contains a positive integer, by negating a nonzero element if necessary. Choose the least positive one, $d$. For any member $a$, divide $a=qd+r$ with $0\leq r<d$. Closure under subtraction puts $r$ in the subgroup, and minimality forces $r=0$. Thus the subgroup is contained in $d\Z$; closure under multiples proves the reverse inclusion. The least positive element of $d\Z$ is $d$, proving uniqueness for nonzero subgroups. A nonzero subgroup cannot equal $0\Z$. Finally $k\mapsto dk$ is an additive isomorphism $\Z\to d\Z$ when $d>0$.
\end{proof}

\KEntry{K07.02.cyclic_abelian_and_subgroups}{Subgroups and images of cyclic groups}
\begin{proposition}\label{prop:cyclic-subgroups-images}
Cyclic groups are abelian. Every subgroup and every homomorphic image of a cyclic group is cyclic. A nontrivial subgroup of an infinite cyclic group is infinite cyclic.
\end{proposition}
\begin{proof}
Two powers of one element commute, since $x^{a+b}=x^{b+a}$. For $H\leq\langle x\rangle$, pull $H$ back along $\Z\to\langle x\rangle$, $k\mapsto x^k$. This preimage is an additive subgroup, hence $d\Z$ by Proposition~\ref{prop:integer-subgroups}. Its image is exactly $H$, so $H=\langle x^d\rangle$. In the infinite case a nontrivial image has $d>0$, and $x^d$ has infinite order by Proposition~\ref{prop:order-divisibility}. Finally $f(x^k)=f(x)^k$ shows that the image is generated by $f(x)$.
\end{proof}

\KEntry{K07.02.order_of_power}{The order of a power}
\begin{proposition}\label{prop:order-power}
For $\operatorname{ord}(x)=n<\infty$ and $k\in\Z$,
\[
 \operatorname{ord}(x^k)=\frac{n}{\gcd(n,|k|)}.
\]
For infinite-order $x$, its zeroth power has order one, and $x^k$ has infinite order when $k\ne0$.
\end{proposition}
\begin{proof}
Let $d=\gcd(n,|k|)$, write $n=dn'$ and $k=dk'$ with $\gcd(n',|k'|)=1$. The least positive $t$ satisfying $(x^k)^t=1$ is the least positive $t$ with $n\mid kt$, or $n'\mid k't$. Bezout's identity for $n',k'$ makes this equivalent to $n'\mid t$. Hence the least such $t$ is $n'$. This includes $k=0$, when $d=n$ and $n'=1$. For infinite order and $k\ne0$, $(x^k)^t=1$ with $t>0$ would give a nonzero annihilating exponent $kt$, contradicting Proposition~\ref{prop:order-divisibility}. The case $k=0$ is the identity.
\end{proof}

\KEntry{K07.02.finite_cyclic_subgroup_classification}{Subgroups and generators in the finite cyclic case}
\begin{theorem}\label{thm:finite-cyclic-subgroups}
If $G=\langle x\rangle$ has order $n\geq1$, then for each positive divisor $d\mid n$ its unique subgroup of order $d$ is $\langle x^{n/d}\rangle$. An element $x^k$ generates $G$ if and only if $\gcd(k,n)=1$, so $G$ has $\varphi(n)$ generators.
\end{theorem}
\begin{proof}
For any subgroup $H$, its inverse image in $\Z$ is $e\Z$ for $e\geq0$, as above. It contains $n\Z$, so $e>0$ and $e\mid n$. Its image is $\langle x^e\rangle$, of order $n/e$ by Proposition~\ref{prop:order-power}. Thus cardinality $d$ forces $e=n/d$, proving uniqueness, and that power has the required order, proving existence. Again by the power formula $x^k$ has order $n$ exactly when $\gcd(k,n)=1$; it generates the whole group exactly in that case, by Theorem~\ref{thm:cyclic-models}. Counting relatively prime residue classes gives the totient count. For $n=1$ there is one element, and $\varphi(1)=1$ under this counting convention.
\end{proof}

\KEntry{K07.02.cyclic_automorphisms}{Automorphisms of cyclic groups}
\begin{proposition}\label{prop:cyclic-automorphisms}
For $G=\langle x\rangle$ of finite order $n$, every automorphism has the form $x\mapsto x^a$ with $\gcd(a,n)=1$, and
$\Aut(G)\cong(\Z/n\Z)^\times$. The automorphisms of $(\Z,+)$ are $k\mapsto k$ and $k\mapsto-k$.
\end{proposition}
\begin{proof}
A homomorphism is determined by $x$, and the rule $x^k\mapsto x^{ak}$ is well-defined modulo $n$ and multiplicative. It is surjective exactly when $x^a$ generates, equivalently when $a$ is a unit modulo $n$ by Theorem~\ref{thm:finite-cyclic-subgroups}. Composition multiplies the exponents: $x^k\mapsto x^{abk}$. Distinct residue exponents give distinct values at $x$, proving that this correspondence is an isomorphism, including $n=1$. An endomorphism of $\Z$ is $k\mapsto ak$, since its value at $1$ determines it. Its image is $a\Z$, which equals $\Z$ exactly when $a=1$ or $a=-1$ by Proposition~\ref{prop:integer-subgroups}.
\end{proof}

\KEntry{K07.02.permutation_convention}{The permutation group and composition}
For any set $X$, let $S_X$ be its bijections, with
$(\sigma\tau)(x)=\sigma(\tau(x))$. The rightmost function acts first. The group laws were proved for bijections in Chapter~\ref{chap:M07.01}. When $X=\{1,\ldots,n\}$ write $S_n$; $S_0$ is allowed. A cycle $(a_1\ a_2\ \cdots\ a_r)$ sends $a_i$ to $a_{i+1}$, sends $a_r$ to $a_1$, and fixes other points. We omit one-cycles when recording a cycle decomposition, but count their points when computing the size of $X$.

\KEntry{K07.02.symmetric_order}{Counting and relabelling permutations}
\begin{proposition}\label{prop:symmetric-order}
For finite $X$ with $|X|=n$, $|S_X|=n!$. A bijection $b:X\to Y$ induces an isomorphism $S_X\to S_Y$, $\sigma\mapsto b\sigma b^{-1}$.
\end{proposition}
\begin{proof}
Order the elements of $X$. There are $n$ choices for the first image, $n-1$ for the next, and successively one for the last, each ordered choice giving exactly one bijection. For $n=0$ the empty choice gives one permutation, in agreement with $0!=1$. Conjugation by $b$ preserves products because $b^{-1}b=\id_X$, and conjugation by $b^{-1}$ is its inverse. Hence it is a group isomorphism.
\end{proof}

\KEntry{K07.02.cycle_decomposition}{Disjoint cycles}
\begin{theorem}\label{thm:disjoint-cycles}
A permutation of a finite set is a product of pairwise disjoint nontrivial cycles, unique up to ordering of the cycles and cyclic rotation of their entries. Disjoint cycles commute.
\end{theorem}
\begin{proof}
Starting with $a$, the sequence $a,\sigma a,\sigma^2a,\ldots$ repeats because $X$ is finite. If $\sigma^i a=\sigma^j a$ with $i<j$, apply $\sigma^{-i}$ to obtain $a=\sigma^{j-i}a$. Let $r$ be the least positive return time. Then these first $r$ points are distinct and $\sigma$ acts on them as one cycle. Two such strings that share a point contain exactly the same points: if $\sigma^i a=\sigma^j b$, every iterate of either is an iterate of the other, allowing negative iterates since the permutation is invertible. Thus these strings partition $X$. On each part $\sigma$ is the cycle just described, so their product agrees with $\sigma$ at every point. Singleton parts contribute identity maps and are omitted. These parts and their cyclic order are determined by $\sigma$ itself, proving uniqueness. Permutations with disjoint supports commute because at a point in either support the other permutation fixes both that point and its image; outside both supports each is the identity.
\end{proof}

\KEntry{K07.02.permutation_order}{Order from cycle lengths}
\begin{proposition}\label{prop:permutation-order}
The order of a finite permutation is the least common multiple of its disjoint cycle lengths, with the empty least common multiple equal to one.
\end{proposition}
\begin{proof}
On a cycle of length $r$, the $k$th power is identity exactly when $r\mid k$, by following a point around the cycle. Since the cycles partition the set, $\sigma^k$ is identity exactly when every length divides $k$. The least positive such integer is their least common multiple. The identity has no nontrivial cycles and already has order one.
\end{proof}

\KEntry{K07.02.transpositions_parity}{Transpositions and sign}
\begin{theorem}\label{thm:permutation-sign}
Every finite permutation is a product of transpositions. The parity of the number of transpositions is independent of the factorization. The resulting sign $\sgn:S_X\to\{1,-1\}$ is a homomorphism; a nontrivial transposition has sign $-1$ and an $r$-cycle has sign $(-1)^{r-1}$.
\end{theorem}
\begin{proof}
Under our composition convention,
\[
 (a_1\ a_2\ \cdots\ a_r)
 =(a_1\ a_r)(a_1\ a_{r-1})\cdots(a_1\ a_2).
\]
Checking the image of each $a_i$ proves this identity; the disjoint-cycle theorem therefore gives a transposition factorization for every permutation. Choose an ordering of $X$ and form its permutation matrix over $\Q$, sending $e_x$ to $e_{\sigma x}$. Composition multiplies these matrices. Define $\sgn(\sigma)$ as its determinant, using the foundational determinant laws of D05. A transposition exchanges two basis columns, so its determinant is $-1$. Multiplicativity of the determinant shows that a product of $k$ transpositions has determinant $(-1)^k$. This both restricts the value to $\{1,-1\}$ and proves parity independent of all choices of factorization. Determinant multiplicativity proves the homomorphism law, and the displayed factorization proves the cycle formula. Relabelling conjugates the permutation matrix, preserving its determinant, so sign is independent of the original ordering. For the empty set the determinant of the empty identity matrix is one. Keeping values in integer signs, rather than reducing into a characteristic-two field, preserves the parity distinction.
\end{proof}

\KEntry{K07.02.alternating_groups}{Alternating groups}
\begin{theorem}\label{thm:alternating-basic}
The even permutations form $A_X=\ker\sgn$. For $n=|X|\geq2$ this is a normal subgroup of index two, of order $n!/2$. For $n\geq3$ it is generated by three-cycles. For $n=0,1$ it is trivial.
\end{theorem}
\begin{proof}
The kernel is normal by Proposition~\ref{prop:kernel-image}. For $n\geq2$ choose a transposition $t$. Multiplication by $t$ bijects the even permutations with the odd permutations, since sign changes by $-1$. The two sets therefore have the same size and partition $S_X$, giving order $n!/2$ and two cosets. A three-cycle is even by Theorem~\ref{thm:permutation-sign}. Conversely write an even permutation as an even number of transpositions and pair consecutive transpositions. Equal transpositions cancel; two sharing exactly one point have product a three-cycle, after naming their three points in their order of action. Two disjoint transpositions satisfy the explicit identity
\[
 (a\ b)(c\ d)=(a\ c\ b)(a\ c\ d),
\]
verified by their action on $a,b,c,d$. Thus each pair is a product of three-cycles, proving generation. When $n=0,1$ there is only the identity permutation.
\end{proof}

\KEntry{K07.02.conjugacy_cycle_type}{Conjugacy and cycle type}
\begin{proposition}\label{prop:permutation-conjugacy}
Two permutations of the same finite set are conjugate in $S_X$ if and only if they have the same multiset of nontrivial cycle lengths. Conjugation relabels the entries of every cycle.
\end{proposition}
\begin{proof}
For any bijection $b$, evaluating on $b(a_i)$ shows
$b(a_1\cdots a_r)b^{-1}=(b(a_1)\cdots b(a_r))$. This proves necessity by uniqueness of disjoint cycles. Conversely pair the cycles of equal lengths and map the successive entries of each first cycle to the successive entries of its partner. The unmatched points are fixed points, and their numbers agree because both permutations act on the same set of cardinality $|X|$ and the sums of the recorded lengths agree. Extend the cycle maps by a bijection between those fixed points. The resulting bijection conjugates the first permutation to the second. This assertion concerns conjugacy in the full symmetric group; it does not assert that conjugacy classes remain unsplit in an alternating subgroup.
\end{proof}

\KEntry{K07.02.cayley}{Cayley's embedding}
\begin{theorem}\label{thm:cayley}
Every group embeds in the permutation group of its underlying set. A finite group of order $n$ embeds in $S_n$.
\end{theorem}
\begin{proof}
For $g\in G$ let $L_g(x)=gx$. Its inverse is $L_{g^{-1}}$, and associativity gives $L_gL_h=L_{gh}$. Thus $g\mapsto L_g$ is a homomorphism to $S_G$. If $L_g=L_h$, evaluate at $1$ to obtain $g=h$. This proves injectivity for finite or infinite groups. In the finite case choose a bijection from the underlying set to an $n$-element set and use Proposition~\ref{prop:symmetric-order} to relabel the image.
\end{proof}
